\chapter{American Options and Early Exercise}\label{ch:dv:american-options-and-early-exercise}

The share trades at 100.30 and goes ex-dividend tomorrow morning, when it
will drop by about the dividend of 1.00. A deep in-the-money call struck at
80, expiring in a month, is on the book of thousands of holders. Exercised
tonight, it delivers a share worth 100.30 with the dividend attached: 20.30.
Held through the ex-date, it is a call on a share worth 99.30, and even with
its month of remaining optionality it is worth about 19.56. The difference, 0.74
a share, goes to whoever is on the other side of the calls that are not
exercised: the writers, who are assigned only by the holders who do exercise.
Every quarter, before every dividend, the same decision is taken or forgotten
across the listed market. This chapter works out when early exercise is optimal,
what the boundary between exercising and holding looks like, how dividends and
rates move it, how to price American options without a full tree, and what the
decision looks like on a real expiry night.

\section{Why and when to exercise early}

An American option can be exercised at any time before expiry (One Quant Book 1,
chapter 23). Its holder chooses a stopping time; the option's value is the
largest discounted expected payoff over all such choices, the optimal stopping
problem of One Quant Book 4, chapter 10:
\[
V_0=\sup_{\tau\le T}\E^{\mathbb Q}\bigl[e^{-r\tau}g(S_\tau)\bigr].
\]
On the tree of chapter 2 this supremum is the backward induction that takes, at
each node, the larger of exercising and holding.

\begin{definition}[Early-exercise premium]\label{def:dv:american-options-and-early-exercise:premium}
The \emph{early-exercise premium}\index{early-exercise premium} of an American option
is the difference between its value and the value of the European option with the
same strike and expiry. It is never negative, and it is zero when early exercise is
never optimal.
\end{definition}

\begin{proposition}[When exercise can be optimal]\label{prop:dv:american-options-and-early-exercise:when}
With $r\ge0$: an American call on a share that pays no dividend before expiry is never
exercised early, and its premium is zero; a call on a share that pays a dividend may be
exercised only just before an ex-date; an American put may be exercised at any time
when rates are positive, and its premium is strictly positive.
\end{proposition}
\begin{proof}
Call: $C\ge S-Ke^{-r(T-t)}\ge S-K$, so the call is worth at least its exercise value
alive (chapter 2). Between ex-dates the share has no dividend, so the argument holds up
to the moment before each ex-date, where the share is about to lose $D$ and the bound
becomes $S-D-Ke^{-r(T-t)}$. Put: deep in the money the put is worth at most $Ke^{-r(T-t)}-S$
as a European, less than $K-S$ when $r>0$; exercising earns the interest on $K$.
\end{proof}

The call and the put give up different things by being exercised: the call gives up the
interest on the strike it pays early and the insurance of its optionality, and gains only
dividends; the put gives up its optionality and gains the interest on the strike it
receives early. \Cref{fig:dv:american-options-and-early-exercise:premium} shows the
put's premium growing with the rate: zero at a zero rate, 0.52 at 5\%, 1.06 at 10\% for
a one-year at-the-money put at 20 volatility.

\begin{omfigure}
\begin{tikzpicture}
\begin{axis}[width=10.5cm,height=4.8cm,
  xlabel={interest rate (\%)},ylabel={early-exercise premium},
  tick label style={font=\small},label style={font=\small},
  xmin=0,xmax=10,ymin=0,ymax=1.15,grid=major,grid style={gray!25},
  table/col sep=comma]
\addplot[omThm,very thick,mark=*,mark size=1.5pt] table[x=rate,y=premium]{figdata/derivatives/06-american-options-and-early-exercise/premium.csv};
\end{axis}
\end{tikzpicture}

\omcaption{\omterm{def:dv:american-options-and-early-exercise:premium}{Early-exercise premium} of a one-year at-the-money American put
(spot 100, volatility 20\%, no dividend) by interest rate: the value of receiving the
strike early. At zero rate there is nothing to gain from exercising. Data: the
tutorial.}\label{fig:dv:american-options-and-early-exercise:premium}
\end{omfigure}

\section{The exercise boundary}

\begin{definition}[Exercise boundary]\label{def:dv:american-options-and-early-exercise:boundary}
The \emph{exercise boundary}\index{exercise boundary} of an American option is the curve
$t\mapsto S^*(t)$ that separates the spots where exercising is optimal from those where
holding is: for a put, exercise is optimal when $S_t\le S^*(t)$. On the boundary the
option equals its exercise value and, in the \omterm{def:dv:black-scholes-three-ways:model}{Black--Scholes model}, meets it with the same
slope (smooth pasting, One Quant Book 4, chapter 10).
\end{definition}

For a put without dividends the boundary rises towards the strike as expiry approaches
(when $r>0$), and as the time to expiry grows it falls towards the boundary of the
perpetual put. That limit has a closed form.

\begin{proposition}[Perpetual American put]\label{prop:dv:american-options-and-early-exercise:perpetual}
In the \omterm{def:dv:black-scholes-three-ways:model}{Black--Scholes model} with $r>0$ and no dividend, a put that never expires is worth
$(K-S^*)(S/S^*)^{-\gamma}$ for $S>S^*$ and $K-S$ below, with $\gamma=2r/\sigma^2$ and
$S^*=K\gamma/(1+\gamma)$.
\end{proposition}
\begin{proof}
With no time dependence the pricing equation is $\tfrac12\sigma^2S^2V^{\prime\prime}+rSV^{\prime}-rV=0$,
solved by $S^{-\gamma}$ and $S$; a put that vanishes at infinity keeps only $AS^{-\gamma}$.
Value matching $AS^{*-\gamma}=K-S^*$ and smooth pasting $-\gamma AS^{*-\gamma-1}=-1$ give
$S^*=K\gamma/(1+\gamma)$.
\end{proof}

With $r=5\%$ and $\sigma=20\%$, $\gamma=2.5$ and the perpetual boundary is 71.43; the
boundary of a one-year put is 81.11 today, and of a three-year put about 76.2
(\cref{fig:dv:american-options-and-early-exercise:boundary}).

\begin{omfigure}
\begin{tikzpicture}
\begin{axis}[width=10.5cm,height=5cm,
  xlabel={time to expiry (years)},ylabel={spot},
  tick label style={font=\small},label style={font=\small},
  xmin=0,xmax=3,ymin=65,ymax=104,grid=major,grid style={gray!25},
  legend columns=3,legend style={font=\footnotesize,at={(0.5,-0.3)},anchor=north,draw=none,fill=none,/tikz/every even column/.append style={column sep=8pt}},
  table/col sep=comma]
\addplot[name path=bd,omThm,very thick] table[x=tau,y=boundary]{figdata/derivatives/06-american-options-and-early-exercise/boundary.csv};
\path[name path=axisbottom] (axis cs:0,65) -- (axis cs:3,65);
\addplot[omThm!15] fill between[of=bd and axisbottom];
\addplot[omDef,thick,dashed] table[x=tau,y=perpetual]{figdata/derivatives/06-american-options-and-early-exercise/boundary.csv};
\addplot[gray,thick,dotted] table[x=tau,y=strike]{figdata/derivatives/06-american-options-and-early-exercise/boundary.csv};
\node[font=\footnotesize,text=omThm] at (axis cs:2.2,68) {exercise};
\node[font=\footnotesize] at (axis cs:2.2,92) {hold};
\legend{exercise boundary,,perpetual limit 71.43,strike}
\end{axis}
\end{tikzpicture}

\omcaption{\omterm{def:dv:american-options-and-early-exercise:boundary}{Exercise boundary} of an American put (strike 100, $r=5\%$, $\sigma=20\%$) by
time to expiry, from a 3\,000-step tree: below the curve the holder exercises. The
boundary rises to the strike as expiry approaches and falls towards the perpetual
put's 71.43 as the horizon lengthens. The steps are the tree's grid. Data: the
tutorial.}\label{fig:dv:american-options-and-early-exercise:boundary}
\end{omfigure}

\section{Dividends and rates as triggers}

For a call the only reason to exercise early is to capture a dividend, and the only
moment is just before an ex-date. The holder compares two values the night before.

\begin{method}[The call's exercise decision before an ex-date]\label{met:dv:american-options-and-early-exercise:decision}
With spot $S$, strike $K$, dividend $D$ and time $T-t$ left after the ex-date:
\begin{enumerate}
\item exercise value: $S-K$ (the share, with its dividend);
\item holding value: at least the European call on the ex-dividend share,
$C(S-D,K,T-t)$;
\item exercise if the first exceeds the second. Since $C(S-D,K,T-t)\ge S-D-Ke^{-r(T-t)}$,
exercise requires $D>K\bigl(1-e^{-r(T-t)}\bigr)$: the dividend must exceed the interest
on the strike over the remaining life. For a deep in-the-money call, whose time value is
little more than that interest, the threshold is close to it.
\end{enumerate}
\end{method}

\begin{example}[The night before the ex-date]\label{ex:dv:american-options-and-early-exercise:night}
Spot 100.30, strike 80, 29 days left, $r=4\%$, $\sigma=25\%$, dividend 1.00. Exercising is
worth 20.30, holding 19.56. The interest on the strike over 29 days is 0.254, and the
dividend above which exercising wins is 0.255: the put component of the call's time
value is worth a tenth of a cent. With a dividend of 1.00 the holder who does not
exercise gives up 0.74 a share (\cref{fig:dv:american-options-and-early-exercise:decision}).
\end{example}

\begin{omfigure}
\begin{tikzpicture}
\begin{axis}[width=10.5cm,height=4.8cm,
  xlabel={dividend paid tomorrow},ylabel={value of the 80 call tonight},
  tick label style={font=\small},label style={font=\small},
  xmin=0,xmax=1,ymin=19.2,ymax=20.6,grid=major,grid style={gray!25},
  legend columns=2,legend style={font=\footnotesize,at={(0.5,-0.3)},anchor=north,draw=none,fill=none,/tikz/every even column/.append style={column sep=8pt}},
  table/col sep=comma]
\addplot[omThm,very thick] table[x=div,y=exercise]{figdata/derivatives/06-american-options-and-early-exercise/decision.csv};
\addplot[omDef,very thick] table[x=div,y=hold]{figdata/derivatives/06-american-options-and-early-exercise/decision.csv};
\draw[gray,dashed] (axis cs:0.255,19.2) -- (axis cs:0.255,20.6);
\node[font=\footnotesize,anchor=west] at (axis cs:0.27,19.35) {threshold 0.255};
\legend{exercise tonight,hold through the ex-date}
\end{axis}
\end{tikzpicture}

\omcaption{The call of \cref{ex:dv:american-options-and-early-exercise:night} the night
before the ex-date: its exercise value does not depend on the dividend, its value held
through the ex-date falls one for one with it. Above a dividend of 0.255 the holder
should exercise. Data: the tutorial.}\label{fig:dv:american-options-and-early-exercise:decision}
\end{omfigure}

Puts move the other way around dividends: a dividend lowers the share after the ex-date
and so raises the put, which makes holding a put through an ex-date more attractive and
early exercise just before it less so. High rates and deep moneyness make early exercise of
puts more likely; high volatility makes it less likely, since it raises the value of
waiting, and so does a dividend yield or a borrow fee (chapter 5), which lowers the forward
the put is waiting for. The same borrow fee works the other way on calls: a call holder who
exercises owns a share that can be lent at the fee, and on a hard-to-borrow share that can
justify exercising a call early with no dividend in sight.

\section{Approximations and bounds}

A tree prices one American option in milliseconds. A market maker who reprices tens of
thousands on every tick needs something faster, and two closed-form approximations are
standard.

\begin{definition}[Quadratic approximation]\label{def:dv:american-options-and-early-exercise:baw}
The \emph{quadratic approximation}\index{quadratic approximation} of Barone-Adesi and Whaley
writes the American value as the European value plus a premium of the form $A(S/S^*)^{q}$,
where the exponent $q$ solves the quadratic that the time-independent part of the pricing
equation imposes, and the critical spot $S^*$ is found numerically from value matching
and smooth pasting at $S^*$. Beyond $S^*$ the option is worth its exercise value.
\end{definition}

The approximation is exact for the perpetual option and for very short expiries, and good
in between (see the table below). The Bjerksund--Stensland approximation is a different
idea: it prices exactly the strategy that exercises at a flat trigger, which is feasible
but not optimal, and so gives a lower bound that is close to the true value.

\begin{center}
\footnotesize
\begin{tabular}{@{}rrrrrrrr@{}}
\toprule
spot & expiry & rate & yield & European & tree (2\,001) & quadratic & error \\
\midrule
100 & 1 & 5\% & 0 & 5.5735 & 6.0902 & 6.0976 & $+0.0074$ \\
90 & 0.5 & 8\% & 0 & 11.2671 & 12.1792 & 12.1151 & $-0.0641$ \\
110 & 0.25 & 3\% & 0 & 1.5070 & 1.5201 & 1.5229 & $+0.0028$ \\
100 & 1 & 5\% & 4\% & 9.0396 & 9.2168 & 9.2444 & $+0.0277$ \\
\bottomrule
\end{tabular}

\smallskip
American puts struck at 100 (volatility 20\%, 30\%, 25\%, 25\% by row). The
\omterm{def:dv:american-options-and-early-exercise:baw}{quadratic approximation} is within three cents except deep in the money at a high rate.
\end{center}

\begin{definition}[Bermudan exercise]\label{def:dv:american-options-and-early-exercise:bermudan}
An option with \emph{Bermudan exercise}\index{Bermudan exercise} can be exercised only on
a finite set of dates before expiry. Its value lies between the European and the American
values and increases with the set of exercise dates.
\end{definition}

\omterm{def:dv:american-options-and-early-exercise:bermudan}{Bermudan exercise} is the rule in callable bonds and in the Bermudan swaptions of One
Quant Book 6, chapter 9, and it is what every numerical method actually prices: a tree
with $n$ steps is a Bermudan with $n$ dates. The value converges to the American one as
the dates multiply, quickly at first: a one-year put exercisable monthly is worth 6.04,
already 90\% of the way from the European 5.57 to the American 6.09
(\cref{fig:dv:american-options-and-early-exercise:bermudan}).

\begin{omfigure}
\begin{tikzpicture}
\begin{semilogxaxis}[width=10.5cm,height=4.8cm,
  xlabel={number of exercise dates (log scale)},ylabel={put value},
  tick label style={font=\small},label style={font=\small},
  xmin=1,xmax=200,ymin=5.5,ymax=6.15,grid=major,grid style={gray!25},
  legend columns=3,legend style={font=\footnotesize,at={(0.5,-0.3)},anchor=north,draw=none,fill=none,/tikz/every even column/.append style={column sep=8pt}},
  table/col sep=comma]
\addplot[omThm,very thick,mark=*,mark size=1.5pt] table[x=dates,y=bermudan]{figdata/derivatives/06-american-options-and-early-exercise/bermudan.csv};
\addplot[omDef,thick,dashed] table[x=dates,y=european]{figdata/derivatives/06-american-options-and-early-exercise/bermudan.csv};
\addplot[omProp,thick,dotted] coordinates {(1,6.0902) (200,6.0902)};
\legend{Bermudan,European 5.5735,American 6.0902}
\end{semilogxaxis}
\end{tikzpicture}

\omcaption{A one-year at-the-money put ($r=5\%$, $\sigma=20\%$) exercisable on 1 to 200
equally spaced dates, the last at expiry. One date is the European option; monthly
exercise captures 90\% of the \omterm{def:dv:american-options-and-early-exercise:premium}{early-exercise premium}. Data: the
tutorial.}\label{fig:dv:american-options-and-early-exercise:bermudan}
\end{omfigure}

\section{The practical exercise decision}

The theory says when; the market decides whether. Exercise is an instruction sent by the
holder, through the broker, to the clearing house before a cut-off; only options at least
one cent in the money at expiry are exercised automatically (One Quant Book 1, chapter 23),
and nothing is automatic before expiry. A holder who forgets to exercise a call the night
before an ex-date loses the difference computed above, and the clearing house's random
assignment passes it to the writers who are not assigned.

\begin{dated}{2026-09}{Calls left unexercised before ex-dates}\label{dat:dv:american-options-and-early-exercise:pool}
A study of US equity call options on shares paying quarterly dividends, January 1996 to
April 2006, found that more than half of the long positions that should have been
exercised on the day before an ex-dividend date were not, costing holders over
USD~491 million over the period; market makers captured most of it through a dividend
spread strategy (chapter 26). The one-cent automatic-exercise threshold at expiry is the
options clearing house's rule for equity and index options.
\end{dated}

Three practical points follow. Desks value their short American calls assuming that
holders exercise optimally, and book the gain when they do not; marking the writer's
position with an assumed failure rate is a model risk (chapter 27). Implied volatilities
of American options are computed with an American pricer or after removing the
\omterm{def:dv:american-options-and-early-exercise:premium}{early-exercise premium}: inverting the European formula at the price of a deep in-the-money
American put gives a volatility that is too high. And near-expiry decisions are made with
the dividend and the borrow of the day, which move the thresholds as much as the spot does.

\section{Tutorial: boundaries, approximations and a decision}\label{tut:dv:american-options-and-early-exercise:tut}

\begin{tutorial}
\textbf{Goal.} Compute the put's \omterm{def:dv:american-options-and-early-exercise:boundary}{exercise boundary} on a tree, compare the \omterm{def:dv:american-options-and-early-exercise:baw}{quadratic
approximation} with the tree, value \omterm{def:dv:american-options-and-early-exercise:bermudan}{Bermudan exercise}, and solve the call's decision
before an ex-date. \textbf{End state:} the four figures of the chapter, the table and the
numbers of \cref{ex:dv:american-options-and-early-exercise:night}.
\begin{tutsteps}
\item \textbf{The boundary}: at each step of the tree, the highest node where exercise
wins.
\omcode{code/firm/american/firm_american.py}{63}{84}{The put's exercise boundary read off a Cox--Ross--Rubinstein tree.}\label{lst:dv:american-options-and-early-exercise:boundary}
\item \textbf{The decision} and its dividend threshold:
\omcode{code/firm/american/firm_american.py}{106}{125}{Exercise or hold the night before an ex-date, and the threshold dividend.}\label{lst:dv:american-options-and-early-exercise:decision}
\item \textbf{Run} \texttt{dv\_american.results()}, \texttt{accuracy\_table()},
\texttt{hook()} and \texttt{fig\_american.py}.
\end{tutsteps}
\textbf{What to change next.} Add a second dividend two weeks after the first and find the
dividend thresholds before each ex-date; then compute the put's boundary with a 4\%
dividend yield and compare with the call's.
\end{tutorial}

\section{Build: the American pricer and exercise helper}\label{bld:dv:american-options-and-early-exercise:american}

\begin{build}
\bfield{Purpose} The miniature firm's listed-options desk prices its American options
with the approximation in the quoting loop and with the tree for risk; on every ex-date
eve it lists the calls, long and short, whose exercise is optimal.
\bfield{Interface} \texttt{baw(spot, strike, t, r, q, vol, right) -> (price, critical)};
\texttt{put\_boundary(strike, t, r, q, vol, n)}; \texttt{bermudan\_put(\ldots, n\_ex)};
\texttt{exercise\_decision(spot, strike, dividend, t\_left, r, vol)};
\texttt{dividend\_threshold(\ldots)}; the tree itself is \texttt{firm.binomial}.
\bfield{Rules} The American call without dividends equals the European; the approximation
returns the exercise value beyond its critical spot; the holding value in the decision is
the European call on the ex-dividend share.
\bfield{Acceptance tests} \texttt{code/firm/american/tests/}: the approximation within
0.07 of a 2\,001-step tree on the four cases of the table; the perpetual limit of the
boundary; the Bermudan value increasing in the number of dates, between European and
American; the threshold dividend just above the interest on the strike.
\bfield{Stretch} The Bjerksund--Stensland lower bound, and a dividend-aware decision that
values holding with the American pricer when a further dividend follows.
\end{build}

\begin{omsources}
\item R.~C.~Merton, ``Theory of rational option pricing'', \emph{Bell Journal of Economics
and Management Science} 4 (1973) 141--183.
\item G.~Barone-Adesi and R.~E.~Whaley, ``Efficient analytic approximation of American
option values'', \emph{Journal of Finance} 42 (1987) 301--320.
\item P.~Bjerksund and G.~Stensland, ``Closed-form approximation of American options'',
\emph{Scandinavian Journal of Management} 9 (1993) 87--99.
\item V.~K.~Pool, H.~R.~Stoll and R.~E.~Whaley, ``Failure to exercise call options: an
anomaly and a trading game'', \emph{Journal of Financial Markets} 11 (2008).
\end{omsources}

\section{Exercises}

\begin{exercise}[$\star$]\label{exo:dv:american-options-and-early-exercise:1}
A one-year American put struck at 100 on a share at 100 ($r=5\%$, $\sigma=20\%$) is worth
6.0902 and the European 5.5735. Give the \omterm{def:dv:american-options-and-early-exercise:premium}{early-exercise premium}, and its size relative to
the European value.
\end{exercise}

\begin{exercise}[$\star$]\label{exo:dv:american-options-and-early-exercise:2}
Give the boundary and the value at spot 100 of a perpetual American put struck at 100 with
$r=5\%$ and $\sigma=20\%$.
\end{exercise}

\begin{exercise}[$\star$]\label{exo:dv:american-options-and-early-exercise:3}
A call struck at 50 has 60 days left after tomorrow's ex-date; $r=5\%$. Below which dividend
is its early exercise certainly not optimal?
\end{exercise}

\begin{exercise}[$\star\star$]\label{exo:dv:american-options-and-early-exercise:4}
Why does high volatility make early exercise of a put less likely, and a high interest
rate more likely? What does a borrow fee do?
\end{exercise}

\begin{exercise}[$\star\star$]\label{exo:dv:american-options-and-early-exercise:5}
Show that the Bermudan put with one exercise date at expiry is the European put, and that
adding dates cannot lower its value.
\end{exercise}

\begin{exercise}[$\star\star$]\label{exo:dv:american-options-and-early-exercise:6}
From the table, which case does the \omterm{def:dv:american-options-and-early-exercise:baw}{quadratic approximation} price worst, and why there?
\end{exercise}

\begin{exercise}[$\star\star\star$]\label{exo:dv:american-options-and-early-exercise:7}
\emph{Coding.} With \texttt{dividend\_threshold}, find the threshold dividend for the call of
\cref{ex:dv:american-options-and-early-exercise:night} with 90 days left instead of 29, and
compare it with the interest on the strike.
\end{exercise}

\begin{exercise}[$\star\star\star$]\label{exo:dv:american-options-and-early-exercise:8}
\emph{Find the flaw.} ``The American put is quoted at 6.09. Inverting the European formula
at 6.09 gives 21.4\% volatility, against 20\% for the European options. The American
options are rich: sell them.''
\end{exercise}

\section{Problem: The Night Before the Ex-Date}

\begin{problem}[{Weekend problem --- who collects the dividend}]\label{pb:dv:american-options-and-early-exercise:1}
The share of the opening trades at 100.30 the evening before it goes ex-dividend by 1.00.
The 80 call has 29 days left after tomorrow, $r=4\%$, volatility 25\%, and an open interest
of 10\,000 contracts (multiplier 100). Experience suggests that 45\% of holders will not
exercise tonight.

\medskip\noindent\textbf{Part I --- The decision.}
\begin{enumerate}
\item Give the exercise value tonight.
\item Give the value of holding through the ex-date.
\item Should a holder exercise? By how much per share?
\item Give the interest on the strike over the remaining 29 days.
\item Give the threshold dividend, and explain why it is so close to the interest.
\end{enumerate}

\medskip\noindent\textbf{Part II --- Who pays.}
\begin{enumerate}[resume]
\item Give the loss of all holders if nobody exercised.
\item Give the transfer to writers if 45\% of holders fail to exercise.
\item Which writers receive it, and how does random assignment decide?
\item Why can a market maker who is short these calls not count on the transfer?
\item What does a writer who is assigned tomorrow morning hold, and what does it owe?
\end{enumerate}

\medskip\noindent\textbf{Part III --- Prices.}
\begin{enumerate}[resume]
\item Where should the call's bid be tonight, relative to 20.30?
\item Tomorrow, what is the call worth if the share opens at 99.30?
\item A trader buys the share and sells the call tonight at 20.30. What does the trade earn
if the call is not exercised against it?
\item And if it is?
\item How would you size such a trade given the uncertainty on the failure rate?
\end{enumerate}

\medskip\noindent\textbf{Part IV --- Judgement.}
\begin{enumerate}[resume]
\item Why do retail holders fail to exercise?
\item Should a broker exercise for its clients automatically?
\item How should a desk mark its short calls tonight?
\item State the \emph{named result}: the threshold dividend and the value transferred to
writers by the holders who fail to exercise.
\item In one sentence: when is an American call worth more than a European one?
\end{enumerate}
\end{problem}

\section{Interview questions}

\begin{interviewq}[$\star$ \iqroles{trader}]\label{iq:dv:american-options-and-early-exercise:1}
Would you ever exercise an American call early? An American put?
\end{interviewq}

\begin{interviewq}[$\star$ \iqroles{researcher}]\label{iq:dv:american-options-and-early-exercise:2}
Why is an American option's value a supremum over stopping times, and how does a tree
compute it?
\end{interviewq}

\begin{interviewq}[$\star\star$ \iqroles{researcher}]\label{iq:dv:american-options-and-early-exercise:3}
Derive the \omterm{def:dv:american-options-and-early-exercise:boundary}{exercise boundary} of a perpetual American put.
\end{interviewq}

\begin{interviewq}[$\star\star$ \iqroles{trader, risk}]\label{iq:dv:american-options-and-early-exercise:4}
You are short deep in-the-money calls the day before a large dividend. What do you expect,
and how do you hedge?
\end{interviewq}

\begin{interviewq}[$\star\star$ \iqroles{developer}]\label{iq:dv:american-options-and-early-exercise:5}
You must price 50\,000 American options per second. Tree, grid or approximation? What do you
check?
\end{interviewq}

\begin{interviewq}[$\star\star\star$ \iqroles{researcher, trader}]\label{iq:dv:american-options-and-early-exercise:6}
How do you compute an implied volatility from an American option price, and what goes wrong
if you use the European formula?
\end{interviewq}
